\subsection{实验现象与问题分析}

实验一使用 ICMP、TCP SYN 和 UDP 探测发现了 \texttt{192.168.4.61}、\texttt{192.168.4.65} 和 \texttt{192.168.4.69} 三台隐藏主机；Nmap 在资产区发现 5 台存活设备，并识别出服务器 \texttt{192.168.4.105} 开放 \texttt{80/tcp}。主机发现用于确认有响应的设备，端口扫描则用于识别开放端口及服务。

实验二发送 30 个 SYN 时，开启和关闭 SYN Cookies 的两组测试均观察到 30 条 \texttt{SYN\_RECV}，开启时 \texttt{TcpExtSyncookiesSent} 未增加。发送量提高到 500 后，两组最大观察值分别为 193 条（关闭）和 385 条（开启）；HTTP 检查均返回 \texttt{401 Unauthorized}。两轮测试的结果还受报文到达、连接超时和采样时刻影响；本次 500 次测试未记录 Cookie 计数器的前后变化。后续对比可同步记录 Cookie 发送计数、半连接状态和 HTTP 响应，以观察防护机制及服务状态。

实验三通过镜像流量读取 TCP 会话信息并构造 RST。攻击期间，受害者服务器 \texttt{192.168.4.105} 上的匹配连接被重置，正常服务器 \texttt{192.168.4.106} 保持可用；停止监听后，大文件下载恢复。小型 HTTP 响应在 RST 到达前已传输完成，体现了报文时序对攻击效果的影响。

实验四依据镜像报文中的 TCP 序列号和确认号，向既有 Telnet 会话注入数据。受害者服务器随后连接到 Attacker 的监听端口，身份检查显示执行权限为 \texttt{uid=1001(lab)}。实验呈现了 Telnet 会话在缺少加密和完整性保护时面临的数据注入风险，注入命令以原登录账户权限执行。

\subsection{安全防护与改进思考}

资产探测可通过网段划分和访问控制加以限制。面对 SYN 洪泛，可结合 SYN Cookies 配置、半连接状态、Cookie 发送计数及 HTTP 响应，观察防护机制和服务可用性。

RST 注入和会话注入都依赖攻击者观察会话流量并发送伪造报文，因此应保护流量镜像端口和路由设备，并在网络边界过滤来源地址异常的报文。应用层应停用 Telnet，改用 SSH 等加密协议；同时限制服务器不必要的出站连接，缩小反向连接可利用的范围。
